\documentclass[11pt]{article}
\usepackage[T1]{fontenc}
\usepackage{lmodern,amsmath,amssymb,amsthm}
\usepackage[margin=1in]{geometry}
\usepackage[hidelinks]{hyperref}
\newtheorem{theorem}{Theorem}
\newtheorem{lemma}[theorem]{Lemma}
\newcommand{\F}{\mathbb F}
\newcommand{\OO}{\mathcal O}
\newcommand{\Gr}{\operatorname{Gr}}
\newcommand{\Hom}{\operatorname{Hom}}
\newcommand{\tr}{\operatorname{tr}}
\title{Integral Splittings in Split and Unitary Rank-Three Compact Models}
\author{\small \href{mailto:admin@cosmosgeek.org}{admin@cosmosgeek.org}\\[-2pt]
\small\href{mailto:b321014@iiit-bh.ac.in}{b321014@iiit-bh.ac.in}}
\date{Working draft for specialist review}
\begin{document}
\maketitle
\begin{abstract}
We compute integral splittings of two maps in depth-one compact models
for rank-three Siegel induction. At a split place a sum over two-by-six
matrices maps onto the Grassmann incidence image. On a nonsingular
central character its retraction averages over $q+1$ lines. We
determine the compact intertwining matrix and construct integral
preimages at its critical parameter. At an unramified inert place a
Steinberg-weighted boundary operator has two rows on ten parahoric
orbits. Two orbit indicators give a right inverse throughout the
unramified parameter germ. Both calculations retain the lattices
and denominators governing coefficient specialization.
\end{abstract}

\section{The two splitting problems}
The pullback of a rank-three Siegel section to a lower-rank induced
representation involves choices of local sections. An explicit
splitting over the coefficient ring makes the image accessible even
after nonflat specialization. We treat two local configurations in
the framework of Skinner and Urban \cite[\S\S10--11]{SU}.

At a split prime the finite compact source consists of functions on
the ordinary Grassmannian of three-planes. A finite oscillator sum
selects the incidence image of the two-planes. We compute a retraction
onto that image and its action on a finite Jacobi character space.
The Grassmann up/down spectrum is classical. We recall the relevant
eigenvalues and compare with Srinivasan \cite{Srinivasan}. The compact
intertwiner then gives a section family through a prescribed critical
vector.

At an inert prime the compact source consists of functions on maximal
isotropic three-planes. We evaluate two boundary rows on the subspace
invariant under an outer Siegel parabolic and a middle Borel. The
geometric series from the affine Weyl cosets has a moving denominator,
while the minor defining the right inverse is constant. The common
application is lifting a prescribed vector or boundary value
integrally without changing its compact level.
\section{The split compact source}
Let $E,D,F$ denote two-dimensional $k$-vector spaces and put
$V=E\oplus D\oplus F$. Throughout, $q$ is an odd prime power.
Let $\mathcal O$ be a discrete valuation ring of characteristic
zero and residue characteristic $p\nmid q$, containing the values of a
nontrivial additive character $\psi:k\to\mathcal O^\times$.
Write $M_i=\mathcal O[\Gr_i(V)]$ for functions on $i$-planes, with
their delta-function bases. The group $\mathrm{GL}(V)$ acts by
permuting subspaces. Define
\[
 (U_i f)(L)=\sum_{\substack{W\subset L\\\dim W=i}}f(W),
 \qquad D_{i+1}=U_i^{\mathrm t},
 \qquad [j]_q=1+q+\cdots+q^{j-1}.
\]
The transposes use the standard basis pairings. One source for this
incidence map is the space of functions $\Phi$ on two-by-six matrices
over $k$. Its summation over matrices with both rows in $L$ gives
\[
 (\lambda\Phi)(L)
   =\sum_{\substack{X\in\operatorname{Mat}_{2,6}(k)\\
             \operatorname{row}(X)\subset L}}\Phi(X).
\]
Row operations by $\mathrm{GL}_2(k)$ preserve this sum.

\begin{lemma}\label{lem:project}
Suppose $[2]_q,\ldots,[6]_q$ are units in $\mathcal O$.
Then $\operatorname{im}\lambda=\operatorname{im}U_2$, and the
integral projector onto this direct summand with kernel $\ker D_3$ is
\[
             P=U_2(D_3U_2)^{-1}D_3 .
\]
\end{lemma}
\begin{proof}
A rank-two matrix with row space $W$ maps under $\lambda$ on its
delta function to $U_2\mathbf1_{\{W\}}$. For a fixed line in a
three-plane there are $[2]_q$ intermediate two-planes. There are
$[3]_q$ two-planes in the whole three-plane. These counts express
the images of rank-one and rank-zero inputs in $\operatorname{im}U_2$.

For completeness the needed eigenvalues of the inclusion matrices follow
from the Grassmann up/down recurrence (compare \cite{Srinivasan}):
\[
 D_{i+1}U_i-U_{i-1}D_i=([6-i]_q-[i]_q)\operatorname{id}.
\]
On two different $i$-planes both off-diagonal counts equal one when
their intersection has dimension $i-1$, and equal zero otherwise. On
a single plane the counts are $[6-i]_q$ and $[i]_q$. Use $D_0=0$.
First $D_1U_0=[6]_q$. Then $D_2U_1$ has eigenvalues
$[2]_q[5]_q$ on $U_0M_0$ and $q[4]_q$ on $\ker D_1$.
These are units, giving a direct sum
\[
 M_2=U_1U_0M_0\oplus U_1\ker D_1\oplus\ker D_2.
\]
The recurrence at $i=2$ gives the eigenvalues of $D_3U_2$ on these
summands, respectively
\[
 [3]_q[4]_q,\qquad q[2]_q[3]_q,\qquad q^2[2]_q .
\]
Invertibility supplies the stated retraction and kernel.
\end{proof}

\subsection{A nonsingular central character}
Use row coordinates for $V$. Consider the subgroup $H$ of block upper
unitriangular matrices in $\mathrm{GL}(V)$ with blocks of size two.
Its centre $Z$ consists of the $(1,3)$-blocks
$S:E\to F$. Fix bases and an invertible matrix $\beta$ such that
$\chi(S)=\psi(\tr(\beta S))$ is a nonsingular additive character. For
an $H$-module let $M[\chi]$ denote its $\chi$-eigenspace for $Z$.
This eigenspace is cut out by the integral idempotent
\[
        e_\chi=q^{-4}\sum_{S\in Z}\chi(S)^{-1}S .
\]
The external group $K=\mathrm{GL}(D)$ acts through the block diagonal
$(1,g,1)$; it normalizes $H$ and fixes $Z$ pointwise. Both $H$ and $K$
act on every $M_i[\chi]$. Put $S_\chi=M_2[\chi]$.

\begin{theorem}\label{thm:fibre}
Assume only $p\nmid q$. There is an $H\rtimes K$-equivariant
isomorphism
\[
 M_3[\chi]\ \simeq\ S_\chi\otimes_{\mathcal O}
                    \mathcal O[\mathbb P(D)]
\]
under which $U_2:S_\chi\to M_3[\chi]$ sends $x$ to the constant tuple
$(x)_{d\in\mathbb P(D)}$. The action on $S_\chi$ is its inherited
action, and $H$ acts trivially on $\mathbb P(D)$. If $q+1$ is a unit,
there is a unique $H\rtimes K$-equivariant projection onto this diagonal.
It sends
\[
 (x_d)_d\ \longmapsto\
   \left((q+1)^{-1}\sum_{d\in\mathbb P(D)}x_d\right)_{d'\in\mathbb P(D)}.
\]
Under the hypotheses of Lemma~\ref{lem:project}, this is precisely
$P$ on $M_3[\chi]$. If $q+1$ fails to be a unit, no such equivariant
projection exists over $\mathcal O$.
\end{theorem}

\begin{proof}
For a central permutation orbit on $\Gr_i(V)$ the character space
has rank one if the character is trivial on its stabilizer; otherwise
the space is zero. Indeed a character function on an orbit is
determined by one value. An incompatible stabilizer forces that value
to vanish. If the projection of a plane onto $E$ is proper, the central
stabilizer contains all maps vanishing on its projection. Nonsingularity
of the trace pairing ensures that $\chi$ is nontrivial on this
subspace. In particular $M_1[\chi]=0$.

A contributing two-plane is hence a graph $E\to D\oplus F$.
The centre translates the $F$-component freely. Thus $S_\chi$ has
a basis indexed by $Z_0\in\Hom_k(E,D)$, of cardinality $q^4$. The
$(1,2)$-block acts by translations of this index; the $(2,3)$-block
acts by multiplication by the additive characters
$Z_0\mapsto\psi(\pm\tr(\beta Z_0V_0))$ as $V_0:D\to F$ varies.
The common sign depends only on the permutation-action convention.
These are all the additive characters of $\Hom_k(E,D)$ by
nonsingularity. Their averaging idempotents project onto individual
basis indices, since $q$ is a unit. Those projections and the
translations give every matrix unit. In particular the commutant of
$H$ on $S_\chi$ consists of scalars.

A contributing three-plane projects onto $E$ with one-dimensional
kernel $\ell\subset D\oplus F$. If $\ell\subset F$, the maps
$E\to\ell$ lie in its stabilizer and $\chi$ is nontrivial there.
Otherwise $\ell$ is the graph of a map from some line $d\subset D$
to $F$. There are $q^2$ choices of this map for fixed $d$. The centre
acts freely, altering the $F$-part of the lifts of $E$.
Modulo this action the remaining datum for a fixed $\ell$ is a map
$E\to D/d$, with $q^2$ choices. It follows that
$M_3[\chi]=\bigoplus_d S_d$, with each $S_d$ free of rank $q^4$.
The unipotent group $H$ preserves $d$ and $K$ permutes the $d$'s.

Let $A_d$ be the component of $U_2:S_\chi\to S_d$, and let $B_d$
be the restriction of $D_3$ to $S_d$. Their product $B_dA_d$ is
an $H$-endomorphism of $S_\chi$, hence scalar. All these scalars
agree by $K$-equivariance. On $S_\chi$, the recurrence in the proof of
Lemma~\ref{lem:project} gives $D_3U_2=q^2(q+1)\operatorname{id}$
because $D_2S_\chi\subset M_1[\chi]=0$. Adding the $q+1$
scalars and canceling this nonzero integer in characteristic zero
gives $B_dA_d=q^2\operatorname{id}$ over $\mathcal O$. These are
maps between free modules of equal rank. Each $A_d$ is an
isomorphism, with inverse $q^{-2}B_d$. Using the $A_d$ themselves
to identify the summands gives the asserted equivariant tensor
description, with $U_2$ diagonal and $D_3=q^2\sum_d$.

An $H$-equivariant map from the tuple module to $S_\chi$ is
$\sum_d b_dx_d$. The action of $K$ forces all $b_d$ equal. Its
restriction to the diagonal is multiplication by $(q+1)b_d$,
giving both the retraction and the obstruction claimed. The
integral change of coordinates given by the $A_d$ remains valid
after coefficient base change.
\end{proof}

\subsection{The compact intertwiner and its critical preimage}
Here take $q$ an odd prime, $F_0=\mathbb Q_q$, and give
$\operatorname{Mat}_3(\mathbb Z_q)$ additive measure one. The standard
intertwiner for the $(3,3)$-parabolic of $\mathrm{GL}_6(F_0)$ integrates
a section on the row chart $[I,T]$, $T\in\operatorname{Mat}_3(F_0)$
\cite[\S11]{SU}. In the lower-block convention, saturation of the rows
by a left factor of determinant valuation $m$ gives $Y^mq^{-3m}$ in
the integrand. This fixes our parameter $Y$.

A matrix entry on the compact level depends only on the dimension
$r$ of intersection of the input and output three-planes. Its
unnormalized value per individual plane is
\begin{equation}\label{eq:schur}
 c_r(Y)=q^{-9}(Yq^3)^r\prod_{i=1}^{r}\frac{1-q^{-i}}{1-Yq^{3-i}} .
\end{equation}
Fix the output spanned by the last three coordinates. Split each
three-block into $3-r,r$, and write
$T=\left(\begin{smallmatrix}A&B_0\\C&D_0\end{smallmatrix}\right)$.
A coordinate input of intersection $r$ has pivots in the first $3-r$
positions of the first block and the last $r$ of the second block.
Gaussian elimination gives the saturated rows
\[
 \begin{pmatrix}I&-U&Z_0&0\\0&V_0&W&I\end{pmatrix},
 \quad V_0=D_0^{-1},\quad U=B_0D_0^{-1},\quad
 W=D_0^{-1}C,\quad Z_0=A-B_0D_0^{-1}C .
\]
They reduce to the specified coordinate plane precisely when all
four new matrices have entries in $q\mathbb Z_q$. The pivot minor
requires $D_0$ to be invertible. The Jacobian for $dT$ is
$|\det V_0|^{-6}$: inversion accounts for $2r$ in the exponent and
the rectangular blocks for $3-r$ each. The character times the
Jacobian is $(Yq^3)^{v_q(\det V_0)}$. Scaling the nine entries
contributes $q^{-9}$ and the factor $(Yq^3)^r$. The last integral is
\[
 \int_{\operatorname{Mat}_r(\mathbb Z_q)}
 s^{v_q(\det V')}\,dV'=\prod_{i=1}^r\frac{1-q^{-i}}{1-sq^{-i}}.
\]
Separate the common valuation of the first column. After taking out
that power of $q$, an integral change of basis sends its primitive
part to the first unit vector; the residual minor is uniformly
distributed. Induction proves the integral and \eqref{eq:schur}.

Let $\binom{i}{j}_q$ denote the Gaussian binomial. The number of input
planes for a fixed intersection dimension is $q^{(3-r)^2}\binom{3}{r}_q^2$:
choose the intersection, its complementary projection, and a graph
into the quotient of the output by that intersection. The four counts
are $q^9,q^4[3]_q^2,q[3]_q^2,1$. Multiplying them by \eqref{eq:schur}
and summing gives the spherical scalar
\[
 m_0(Y)=\prod_{i=1}^3(1-Yq^{-i})\big/\prod_{j=0}^2(1-Yq^j).
\]
Writing $N=c/m_0$ for the operator fixing constants, its entry at
intersection dimension $r$ is
\[
 d_r(Y)=q^{-9}(Yq^3)^r
 \frac{\prod_{i=1}^r(1-q^{-i})\prod_{j=0}^{2-r}(1-Yq^j)}
      {\prod_{i=1}^3(1-Yq^{-i})}.
\]
These identities hold first for small $|Y|$ and then as rational
identities. Set $Y=X^{-1}$. If $q(q^2-1)(q^3-1)(q^4-1)$ is a
coefficient unit, all remaining denominators are units in the
germ $\mathcal O[[X-q]]$. At $X=q$, the four entries satisfy
$d_0=d_1=0$, $d_2=1/([3]_q[4]_q)$ and $d_3=[3]_qd_2$. The
$(L,L')$ entry of $U_2D_3$ is $\binom{r}{2}_q$, so
\[
 N(q^{-1})=\frac{U_2D_3}{[3]_q[4]_q},\qquad
 g=[3]_q[4]_q U_2(D_3U_2)^{-2}D_3 f \quad(f\in U_2M_2)
\]
satisfies $N(q^{-1})g=f$. The three eigenvalues in
Lemma~\ref{lem:project} are units under the present hypotheses.
Their monic cubic identity, computed in characteristic zero, holds on
the integral module by torsion-freeness; its unit constant coefficient
gives the inverse even when $[5]_q$ or $[6]_q$ is a nonunit.
Thus $N(X^{-1})g$ is an integral family with critical value $f$.

\section{The inert compact source}

Let $q$ be an odd prime, $k=\F_{q^2}$ and
$\bar x=x^q$. On the $k$-space with basis
$e_1,e_2,e_3,f_1,f_2,f_3$ use the split skew-Hermitian form with
$\langle e_i,f_j\rangle=\delta_{ij}$ and
$\langle f_i,e_j\rangle=-\delta_{ij}$. Write
$V_o=\langle e_1,e_2,f_1,f_2\rangle$,
$V_m=\langle e_3,f_3\rangle$,
$F_o=\langle f_1,f_2\rangle$ and $E_m=\langle e_3\rangle$.
The stabilizer of $F_o$ in $U(V_o)$ times the stabilizer of $E_m$ in
$U(V_m)$ acts on the maximal totally isotropic three-planes of
$V_o\oplus V_m$. Let $M$ be the functions constant on its orbits.

There are six decomposed orbits. For a three-plane
$L=L_o\oplus L_m$ the outer index is
$d=\dim(L_o\cap F_o)$, and the middle index records whether
$L_m=E_m$. Use the following labels:
\[
\begin{array}{c|cc}
d&L_m=E_m&L_m\ne E_m\\\hline
2&9&1\\1&7&3\\0&0&5
\end{array}
\]
The three outer positions can be seen by first fixing the intersection
in $F_o$ using its general linear stabilizer. On the remaining
hyperbolic quotient the outer plane is transverse to the remaining
part of $F_o$. Its graph is carried to the standard transverse plane
by a unipotent in the stabilizer of $F_o$. In the middle plane the
stabilizer of $E_m$ is transitive on its complement among isotropic
lines, by the same graph description.

Four further orbits consist of graphs. In fact, the projection of $L$
to $V_o$ or $V_m$ has orthogonal complement $L\cap V_o$ or $L\cap V_m$
respectively, since $L=L^\perp$. If $L\cap V_m$ has dimension one,
the middle projection is that same line and $L$ decomposes. Otherwise
the middle projection is onto and $\ell=L\cap V_o$ has dimension one.
Modulo $\ell$, $L$ is the graph of an anti-isometry
$V_m\to\ell^\perp/\ell$ (the orthogonal here is in $V_o$).
The outer line $\ell$ is either in $F_o$ or outside it. There is
one stabilizer orbit for each position: in the first case use the
Levi action on $F_o$, and in the second first move the projection of
$\ell$ onto the complementary isotropic two-plane by the Levi.
The remaining graph coordinate is removed by the unipotent radical,
since an isotropic graph on a line extends to an isotropic graph on
that two-plane (extend its Hermitian matrix by one row and column).

In each case the stabilizer of $\ell$ induces the full Borel of the
hyperbolic plane $\ell^\perp/\ell$ fixing the image of
$F_o\cap\ell^\perp$. Indeed, at the representatives $\ell=f_2$ and
$\ell=e_2$, the Borel on $\langle e_1,f_1\rangle$ acts while fixing
the other hyperbolic pair. A graph is therefore classified by whether
$E_m$ maps to the distinguished line in that quotient. Write
\[
\begin{array}{c|cc}
 &\text{equal}&\text{different}\\\hline
\ell\subset F_o&8&4\\
\ell\not\subset F_o&6&2 .
\end{array}
\]
Together these tables give a basis $f_0,\ldots,f_9$ of $M$, the
indicator functions of the corresponding orbits, over any coefficient
ring.

\subsection{The boundary operator}

Let $K/\mathbb Q_q$ be unramified quadratic. Use the same split form
and its integral basis over $\OO_K$. For
$h=\left(\begin{smallmatrix}a&b\\ c&d\end{smallmatrix}\right)
\in\mathrm{SL}_2(\mathbb Q_q)$ and $z\in\{f_2,e_2\}$, form the three
rows
\begin{equation}\label{rows}
       d e_3+f_1+c f_3,\qquad e_1+b e_3+a f_3,\qquad z .
\end{equation}
They span a maximal isotropic space: the pairing of the first two
rows is $-1+ad-bc=0$. Intersect their $K$-span with the standard
$\OO_K$-lattice and reduce the saturated lattice modulo $q$. Call
the resulting orbit $j_z(h)$.

We use the affine-Iwahori row normalization in which
\[
 m_r=\operatorname{diag}(q^r,q^{-r}),\quad
 w=\begin{pmatrix}0&-1\\1&0\end{pmatrix},
 \quad r\in\mathbb Z.
\]
These are the translations $m_r$ and reflected translations $m_rw$
for the affine double cosets of $\mathrm{SL}_2(\mathbb Q_q)$.
On the Steinberg Iwahori-fixed line the characteristic Hecke operator,
normalized by division by $H=\operatorname{vol}(I)$, acts by
$(-1)^{\ell(g)}$ for the affine Weyl element $g$: every simple
generator acts by $-1$ and on a reduced word the operators multiply.
In particular the signs for $m_r$ and $m_rw$ are $+$ and $-$.
We keep $H$ in the row formula. The coordinate convention
\eqref{rows} is the row-space restriction used in the project archive
for the local Klingen pullback; the underlying pullback construction
is described in \cite[Section~11]{SU}.

For $X$ an unramified parameter write $u=X/q^3$. Define our two
normalized boundary rows on $f\in M$ by the explicit coset sums
\begin{equation}\label{definition}
 B_z(f)=H\sum_{r\in\mathbb Z}u^{|r|}
       \bigl(f(j_z(m_r))-f(j_z(m_rw))\bigr),
       \qquad z=f_2,e_2 .
\end{equation}
Here $f(j)$ is the value of the function on orbit $j$. Formula
\eqref{definition} can be used directly and fixes all normalization
choices. The factor $u^{|r|}$ uses the lower-block determinant
valuation $-|r|$. Indeed, saturating the rows \eqref{rows} at either
representative for $r\ne0$ multiplies exactly one row by $q^{|r|}$.
In lower-block induction with parameter
$X=\chi(q)^{-1}q^{-2s}$ this supplies
$\chi(q)^{-|r|}q^{-2(s+3/2)|r|}$. Weight two in this convention is
$s=-1/2$, $\chi(q)=1$, hence $X=q$.

\begin{theorem}
Let $t=X/(q^3-X)$. With columns $f_0,\ldots,f_9$ and rows in
the order the vectors $f_2,e_2$ for the third row in \eqref{rows}, the
matrix of $B/H$ is
\[
\begin{pmatrix}
0&t&0&-t&-1&0&0&t&1&-t\\
t&0&-1&t&0&-t&1&-t&0&0
\end{pmatrix}.
\]
If $H$ and $q(q^2-1)$ are units in the coefficient ring, then in
the formal germ at $X=q$ the map has the constant right inverse
\[
 (y_1,y_2)\longmapsto H^{-1}(-y_1 f_4+y_2 f_6).
\]
Its kernel is free of rank eight over that germ and the splitting
commutes with coefficient base change.
\end{theorem}

\begin{proof}
For $m_r$ the first two rows in \eqref{rows} are
$q^{-r}e_3+f_1$ and $e_1+q^r f_3$.
For $m_rw$ they are $f_1+q^{-r}f_3$ and $e_1-q^re_3$.
Use the unit pivots after saturation to read off the following orbits:
\[
\begin{array}{cc|ccc}
z&h&r<0&r=0&r>0\\\hline
 f_2&m_r&1&8&7\\
 f_2&m_rw&9&4&3\\
 e_2&m_r&3&6&0\\
 e_2&m_rw&7&2&5
\end{array}
\]
For example with $r>0$ the first pair reduces to $e_3,e_1$ and
the second to $f_3,e_1$. With $r<0$ they reduce to $f_1,f_3$
and $f_1,e_3$. At $r=0$ the outer intersection is $\langle z\rangle$.
The graph sends $e_3$ to the image of $f_1$ for $m_0$ and to the
image of $e_1$ for $m_0w$, up to a unit sign, giving respectively
the closed and open graph positions in the second orbit table.

Each positive or negative tail has weight
$\sum_{r\ge1}u^r=u/(1-u)=t$, initially for $|u|<1$ and then
rationally. Substitution in \eqref{definition} gives the matrix.
In the formal germ $q^3-X$ has unit constant term. Columns 4 and 6
are $(-H,0)$ and $(0,H)$ for all $X$. If $(b_j,c_j)$ is the $j$th
column of $B/H$, the remaining columns yield the kernel basis
$f_j+b_j f_4-c_j f_6$, $j\ne4,6$. These formulas persist over every
base change of the germ.
\end{proof}

The finite model accounts for two local boundary coordinates. The unit
minor supplies an integral correction in this module, even when the
weight varies. Applying it to an additional toric period, a different
compact level, or a global cusp form requires the maps to those
particular spaces.

\section{Using the splittings}
On the split nonsingular character space take a tuple supported at
one line $d_0\in\mathbb P(D)$, with value $x\in S_\chi$. Projection
gives $x/(q+1)$ in every component. At a coefficient prime dividing
$q+1$, the diagonal embedding still exists over the original lattice
while equivariant retraction fails. At a nonresonant prime the compact
matrix above lifts a prescribed critical coherent section to the
unramified family at the same level.

On the inert source, consider $u_0=f_0+f_6$. At $X=q$ the first
row of $B$ on $u_0$ vanishes; the second has value $Hq^2/(q^2-1)$.
Hence $u_0-q^2 f_6/(q^2-1)$ lies in the critical kernel. In the
moving germ the coefficient of $f_6$ for this correction is
$1+t=q^3/(q^3-X)$. Using the fixed critical coefficient throughout
would generally leave a nonzero boundary. The two constant pivot
columns explain why this correction uses no division by $X-q$.
A functional that factors through these two boundary rows vanishes on
their common kernel. Applying a correction to a different period
requires its own comparison map.

\begin{thebibliography}{9}
\bibitem{Srinivasan} M.~K. Srinivasan,
\emph{The Goldman--Rota identity and the Grassmann scheme},
\href{https://arxiv.org/abs/1304.6244}{arXiv:1304.6244}.
\bibitem{SU} C. Skinner and E. Urban,
\emph{The Iwasawa main conjectures for $\mathrm{GL}_2$},
\href{https://www.math.columbia.edu/~urban/eurp/MC.pdf}{authors' paper},
Sections 10--11, especially Proposition 11.2.3.
\bibitem{record} Project research record of 17 November 2026,
checkpoint 48, \emph{Split finite $\mathrm{GL}_6$: the actual
coherent incidence image and its nonsingular Jacobi projector} and
\emph{Split $\mathrm{GL}_6$: the full compact standard intertwiner};
checkpoint 63, \emph{An actual all-weight minimal-parahoric boundary
right inverse} and \emph{The ten minimal compact orbits: exact native
boundary rows}.
\end{thebibliography}
\end{document}
